\documentclass[12pt]{article}
\usepackage[utf8]{inputenc}
\usepackage[T1]{fontenc}
\usepackage{hyperref}
\title{Yassic Language Resources for K-12 Curriculum Development}
\author{Compiled by Claude Sonnet 5.5 (claude-sonnet-5-5) for Andr\'as Kornai}
\date{}
\begin{document}
\maketitle

\section{Introduction}
Yassic, also called Jassic (ISO 639-3 ysc), is an extinct language of Hungary. It was spoken by the Jász, a group of Alanic origin, and is usually described as a dialect of Ossetian, an East Iranian language. Estimates of when it died out range from the 15th to the 17th century. The only known record is a one-page glossary of about 34 words, found in the Hungarian National Széchényi Library in the 1950s. Several of the pages I found are Wikipedia-derived and flagged for missing citations.

\section{Primary Educational Resources}
None exist. No primers, textbooks or school programmes were found.

\section{Linguistic and Reference Materials}
\begin{itemize}
\item Wikipedia, Jassic dialect: \url{https://en.wikipedia.org/wiki/Jassic_dialect}
\item Wikipedia, Ossetians, for the related living language: \url{https://en.wikipedia.org/wiki/Ossetians}
\end{itemize}

\section{Religious and Community Resources}
\begin{itemize}
\item Wikipedia, Jasz people (history of the community): \url{https://en.wikipedia.org/wiki/Jasz_people}
\end{itemize}

\section{Digital and Audio Resources}
None found.

\section{Curriculum Development Context}
With only a 34-word glossary, a language curriculum is not possible. A heritage curriculum could cover Jász history and identity in Hungary. It could also teach Ossetian, the closest living relative, and the glossary itself. Anything of this kind would need specialist support from Iranian studies scholars and Jász community organisations. I did not find the glossary online or verify the primary sources.
\end{document}